\section{Possible Research}
\label{sec:possibleResearch}

This section contains research ideas that may be persued if time
fermits or if other ideas are unsuccessful.

  Here we describe several ways we will
extend this work.

\begin{figure}
\centering

\includegraphics[width=.2\linewidth]{covar_1}
%\includegraphics[width=.2\linewidth]{covar_2}
\includegraphics[width=.2\linewidth]{covar_3}
%\includegraphics[width=.2\linewidth]{covar_4}
\includegraphics[width=.2\linewidth]{covar_5}
%\includegraphics[width=.2\linewidth]{covar_6}
\includegraphics[width=.2\linewidth]{covar_7}
%\includegraphics[width=.2\linewidth]{covar_8}
%\includegraphics[width=.2\linewidth]{covar_9}
%\includegraphics[width=.2\linewidth]{covar_10}
%\includegraphics[width=.2\linewidth]{covar_11}
%\includegraphics[width=.2\linewidth]{covar_12}
%\includegraphics[width=.2\linewidth]{covar_13}
%\includegraphics[width=.2\linewidth]{covar_14}

\caption{False color images generated from 1, 3, 5, and 7 days
(left--right) of visible satellite images.   The color of each pixel
is determined by the correlation of the time-series of intensities at
that pixel with the time-series of the middle pixel (white means the
correlation is one).  Notice that with sufficient samples the spatial
dependence is roughly isotropic.  It would be possible to make a
reasonable estimate of distance from the middle pixel based only on
the correlation.}

\label{fig:timePixelCorr}

\end{figure}

\subsubsection{Improving Temporal Sampling by Investigating Spatial
Dependence} 

The most critical part of our localization algorithm is the
computation of correlation between videos (or between videos and the
time-series of a satellite pixel).  These correlations depend on the
spatial dependence of the signal being extracted from the webcam.  If
there is no spatial dependence then this method of localization will
not work.  In this objective we will explore the relationship between
spatial dependence, temporal sampling, and localization accuracy with
the goal of improving our temporal sampling scheme.

We believe that localization is most accurate when the spatial
dependence is isotropic, broad, has a distinct peak at zero distance,
and decreases monotonically as distance increases.  Together these
properties mean that we can sort a large range of geographic distance
between sample points using correlation.  This, unfortunately, is not
true for all natural signals.  Initial results from simulation with
visible satellite imagery (each pixel is a proxy for a camera) show
the following:
\begin{itemize}

  \item given sufficient samples the spatial dependence of clouds
    (usually) meets the conditions discussed above,

  \item the spatial dependence is more likely to be anisotropic near
    the ocean shore, and 

  \item at some point, depending on the temporal variability of the
    signal, it becomes useless to increase the temporal sampling rate.

\end{itemize}

\figref{timePixelCorr} shows false-color images that reflect the
correlation between the sensor readings of a single pixel with all
other pixels over different numbers of days. The correlation scores
are higher for nearby pixels.  This is especially true as the number
of days of measurements increases.


We plan on further exploring the interaction between temporal sampling
rate and spatio-temporal properties of the underlying signals to
obtain improved sampling schemes.  The temporal sampling schemes
should avoid over-sampling but should obtain accurate correlation
estimates quickly.   

\subsubsection{Estimating a Depth Map by Relative Localization of Pixels}

Our goal is to create an automatic method for creating a scene depth
map from a sequence of images captured from a static camera.  We will
use the correlation between pixel values over time due to changes
induced by clouds as the primary cue.  This cue is unexplored in the
literature and has significant overlap with our correlation-based
camera localization method.

Recent work on depth estimation for outdoor scenes includes techniques
using optical effects due to particulate
matter~\cite{contrastRestoration, novelDepth2009} and strong
structural priors~\cite{closingTheLoop}. Our approach is more similar
to work on depth-estimation using structured light
patterns~\cite{VaqueroCVPR2009}.  These approaches require a tight
coupling between the camera and the structured light source.  Our
approach differs in only requiring a stochastic description of the
temporal properties of the light pattern (\ie we need to know the
shape of the clouds and how far apart they are). 

Our approach is as follows.  Given a sequence of images, each with $N$
pixels, we compute the correlation, $s_{i,j}$, for each pair of
pixels. We then estimate a mapping (this is a challenging open issue)
that converts the correlation scores into distances in world
coordinates, $d_{i,j} = f(s_{i,j})$. 

A \naive approach to estimating world coordinates from the distance
matrix would be to use multi-dimensional scaling (MDS).  However, MDS
does not use the additional constraints provided by camera geometry
and can therefore create highly improbable depth maps. The remainder
of this section describes our method for including camera calibration
in the process.

We assume the camera calibration is in the form of a matrix, $U \in
\R^{N \times 3}$, where each row represents the unit ray, in world
coordinates, that corresponds to the view direction of a pixel.  We solve for the depth map,
$x^*\in\R^{N\times1}$, by estimating a world position for each pixel
subject to the constraint that it lies along the rays specified by the
camera calibration.  This optimization can be conveniently specified
using an inner-product representation as follows:
\begin{equation}
  x^* = \argmin_x \| R(x)  - M(x) \|_F.
\end{equation}
This formulation minimizes the difference in two inner product
matrices, $R$ and $M$. The ray-based matrix,
\begin{equation}
  R(x) = \diag{x}U\trans{U}\diag{x},
\end{equation}
is the the gram matrix of the world points that correspond to the
known depth map and the known pixel rays.  The image measurement-based
gram matrix, $M$, is defined element-wise as:
\begin{equation}
  M_{i,j}(x) = \frac{1}{2}\left( x_i^2 + x_j^2 + d_{i,j}^2 \right).
\end{equation}
This is a the standard transformation between a distance matrix and a
gram matrix with (conveniently) the assumption that the camera center
is the origin.

We minimize the objective function using a gradient descent method
(\textbf{fminunc} in Matlab).  The formulation has a straightforward
analytical gradient and Hessian computation.  Initial results using
this method with simulated data show promising results.  The next step
is to extend this to real image data.  For real data the challenge
will be estimating the mapping from correlation to distances.

%
%
%
\subsection{Moving Beyond PCA as an Image Model}
\label{sec:generative}

(this section need a complete rewrite)

The foundation of this dissertation is understanding and using the
relationship between signals in the world and image appearance. In the
preliminary work we showed that PCA is an effective means for
extracting image space components that relate to unseen latent
factors.  We also showed that by conditioning the set of images
decomposed by PCA it is possible to extract components that do not
depend on the conditioning variable (\eg PCA on all the noon images
for a week does not depend significantly on sun position).
Unfortunately PCA assumes a very simple, and often incorrect, noise
model.  Consequently, PCA sometimes finds components and coefficient
time-series that are irrelevant to the desired task.

In this objective we will explore alternative generative models.  We
hope to build a simple model that can accurately estimate the
distribution of scene appearances in the future.

\subsubsection{Motivation}

Constructing an accurate generative model of a natural scene is a
challenging task.  In the context of this dissertation, having a
better generative model would improve signal extraction and scene
structure estimation.  It could also help surveillance tasks such as
foreground object detection, object tracking, and anomaly detection.
The AMOS dataset provides a foundation for learning parameters of a
generative scene model.

\subsubsection{Approach}

The first method will consider is a simple extension of our PCA based
model.  We will build a PCA basis as before but we will estimate a
conditional distribution over the PCA coefficients using an
independent gaussian process regression (GPR) model for each
coefficient.  Preliminary experiments with this model have shown
promise in accurately predicting future appearance of the scene.  A
critical component of a GPR model is a similarity function between the
conditioning variables.  Preliminary experiments indicate that
including similarity functions based on the time-of-day give
significant improvements over purely time-based similarity.

However, there are several limitation of the proposed method.  One is
in the use of linear combinations of PCA basis images to represent the
scene.  Another limitation is the use of independent gaussian process
models for each PCA coefficient.  We will investigate alternative
models.

One possible alternative is to more explicityly model physical
effects.  We are considering adding models that consider the
relationship between scene geometry, illumination, particulate matter
in the air, and lighting~\cite{basriPhotometric,sheddingLight}.

\subsubsection{Open Issues}

\begin{itemize}

  \item Could this be used for better signal estimation?
   
  \item How much would it help to simultaneously solve for multiple
    cameras?

  \item Are there better ways than PCA to decompose our scene?

  \item How does this work for foreground object detection? 

  \item In what circumstances is the full generative model better than
    a descriptive model (\eg a prior over the histogram of gradients)?

\end{itemize}


%
%
%
\subsection{Camera Localization Methods for Moving Cameras}
\label{sec:movingLocalization}

The localization method described in \secref{localizationPrelim}
assumes that the camera is static.  The goal of this section is to
remove this limitation to create a method for localizing a moving
camera in an unknown environment using cues from natural environmental
changes.  Making this change requires a new signal extraction methods
and a new location estimation algorithm.

\subsubsection{Motivation}

There are many sources of time-series images captured from moving
cameras.  Localizing such cameras has traditionally been performed in
two ways: frame-frame registration for small image motions and and
frame-model registration for large motions in known environments.  Our
problem statement cannot addressed by either of these methods.

\subsubsection{Approach}

Our approach is to first combine one or more images to obtain a
distribution of possible camera locations.  Then we combine these
distributions using a filtering framework.  We are currently
considering options including the Kalman filter or monte-carlo
localization~\cite{mcl}.

In both techniques we must create a distribution over the geographic
locations of the camera from time-stampped images.  Given that we are
using a moving camera it will be more difficult to extract such a
distribution (using SVD and using correlation for similarity will not
be sufficient).  One approach is to replace SVD and continue to use
correlation.  This would fit naturally into the Kalman filter and
would use the uncertainty estimates~\secref{localizationError}.  The
other approach is to directly estimate the distribution based on a
single image.  For example, if we have a time-stampped image of a
sunny day the location of camera is constrained to be on a known
portion of the globe.  Given sufficient images with these constraints
the location, and the path, of the camera can be estimated.

The feature extraction problem is more difficult with a moving camera.
We expect that estimating the average amount of sunlight will be easy
but that estimating the weather conditions will require more effort.
We will focus on the sunlight cue and explore local-feature based
regression for weather estimation if time permits.  The AMOS dataset
will be useful for learning the regression model.

The MCL based approach has a distinct advantage because it can handle
very general distributions.  In addition to being able to include
broad distribution based on sunlight conditions, it supports sparse
distributions like the distribution given that the image contains a
beach.

\subsubsection{Open Issues}

\begin{itemize}

  \item Is a Kalman filter sufficient?

  \item What is the correct density representation (samples, a
    Gaussian, a fixed grid)?

\end{itemize}


\subsection{Full camera calibration methods}

this has been well played already... do I really want to add
more to this?

Understanding the calibration of the camera can be useful first step
in using images from any camera.  Knowing the calibration reduces the
uncertainty in the real world structure and appearance of the imaged
objects.   Automatic methods for camera calibration are important for
our work because it would be infeasible to manually calibrate all the
cameras in the dataset.  In addition, automatic calibration methods
are of general interest to the field.

\subsubsection{Specific aims}

\subsubsection{Approach}

To determine the remaining camera parameters we will use a number of
cues.  The first cue we will use is the appearance of the sky.  There
exist analytical models of the appearance of the sky for a given sun
position and atmospheric conditions.   We propose to solve the inverse
problem of estimating the camera parameters given a set of appearances
of the sky.

\subsubsection{Progress}

Our work on estimating additional camera parameters appeared in WACV 
2008~\cite{jacobs08geoorient}.  This was followed by work in
estimating the full geometric camera calibration by someone else CITE.


\subsection{Catalog of transient objects}

Thusfar we have focused on building and using models of the
conditional relationship between underlying varaitions and the images.
However, even with precise modeling there will be pixel intensities
that are unexpected. 

\subsubsection{Specific aims} 

We will catalog and analyze the unmodeled intensity variations (e.g.,
snow on the ground, pedestrians, or building construction).  Our goal
is to extract a catalog of the objects, spatially-coherent unmodeled
phenomena, that appear in the scene.

We will create a method of collecting the visual phenomena that are
not due to weather conditions and the sun position.  We will also
create an interface to browse the object catalog for a given camera.

\subsubsection{Approach}

Our approach for detecting the anomalous visual phenomena is to first
model the expected range of appearances using a conditional
(parametric/non-parametric) model.  The image-space residuals from
this conditional model (\ie the low probability pixels) will be
grouped together and added to the catalog.  We expect to experiment
with several methods of modeling the conditional distrubtion on
intensities.

Using this set of cataloged items for a single camera we will provide
a web-based user interface for browsing the object catalog.

Our approach to modelling expected intensities is similar to that used
in the surveillance domain.  The differences is that we use a model
conditioned on signals from the world because we have lower frame-rate
video. 

\subsubsection{Progress}

We are currently working on an initial implementation of the
conditional background model.

\section{Some other Stuff I want to work on too but I probably
don't have time so I should probably not think about it anymore}

%
%
%
\subsection{signals supervised}
\label{sec:signals}

Different latent factors induce different types of variations in
images.  Understanding this relationship is of fundamental importance
to our goal of using a network of webcams to measure the world.

\secref{signalSupervised}

\subsubsection{Motivation}

\subsubsection{Approach}

The challenge with both is that images may vary more due to other
causes than due to the variations related to the signal of interest.
Thus we explore efficient mechanisms to learn features that are
invariant to other scene variations, as well as tools to condition our
feature set on the known causes of image variation.

\subsubsection{Open Issues}

\begin{itemize}

  \item This is not cool if we just do image regression.  I have to
    develop a new technique for image regression and apply it to a
    cool problem.
    
\end{itemize}


%
%
%
\subsection{Characterize Error in Camera Localization from
Geo-temporal Signals}
\label{sec:localizationError}

The camera localization methods described in
\secref{localizationPrelim} gives a single point estimate of the
location of a static camera from a set of images.  We hope to
characterize the causes and the magnitude of errors in the estimate.

The most critical part of our localization algorithm is the
computation of correlation between videos.  This correlation depends
on the underlying spatial-correlation of the signal (\eg degree of
cloudiness, average lighting intensity) that is being
imaged by the webacm.  If this  the underlying signal gives
uniformative camera-camera correlations or the image-image correlation
computation is incorrect (\eg due to insufficient sampling of the
spatial function) then the localization results will be inaccurate.

There is freedom in choosing the underlying spatial function and in
choosing how and when to sample an image.  In this objective we
explore some of these factors with the goal of understanding the
impact they have on localization accuracy.

\subsubsection{Motivation}

Understanding the sources of error will help us improve the accuracy
of our localization algorithm.  It will inform our effort to improve
the sampling strategy used to collect image data.  It will also enable
principled combination of multiple estimates (for an example see
\secref{movingLocalization}) and fusion of other cues, such as the
appearance of recognizable landmarks.

The localization algorithm depends on the relationship between images
and latent factors in the world.  Understanding how uncertainties in
these relationships affect the localization procedure will improve our
understanding of the underlying latent factors.

\subsubsection{Approach}

We hypothesize that there are four primary sources of errors in
computing correlation between cameras: low temporal sampling rate, low
spatial sampling rate, spatial-temporal functions with low
discriminative power (\eg when correlation is due to the length of a
day and the current time is at a solar equinox), and mismatches
between signals that are not modeled by using correlation as a
simliarity metric (\eg non-linear camera response or temporal lag).
We will focus on specific instances of the later two in our
investigation. 

To evaluate the impact of sampling rates we will conduct a number of
psuedo-simulation studies.  We will forego the use of cameras and
instead focus directly on sampling of real spatial-temporal functions.
We will consider pixels in a real geo-registered satellite image and
compute the spatial correlation of pixels using different sampling
rates.  These pixel-pixel correlations will be used as input to our
localization algorithm and the dependency of the errors on the
sampling rates will be evaluted.

To explore the issues of spatial functions with low discriminative
power we will investigate camera-camera correlations based on amount
of sunlight.  We choose this because it is sometimes discriminative,
it is easy to measure from images, and it has a well-known underlying
function (\ie the time when sunrise and sunset on a particular day and
location can be easily computed).  We expect localization accuracy
using sunlight-based correlation to depend on the day of the year and
the location of the camera.

In the the preliminary work we expected simple signal correlation to
depend on the geo-spatial proximity of cameras but there may be better
similarity measures for comparing two cameras.  We will explore one
alternative which is to compute a non-linear mapping from images to
signal values.  The correlation computation does not consider the
impact of the non-linear camera response that is caused by the cameras
automatic exposure compensation.  As an example, consider the case of
two cameras: one with good low-light response and one with poor
low-light response.  Even if they are in the same location the camera
with good low-light response will appear to have a longer day and the
cameras will likely be located differently.

We will compare the ground-truth location and the ground-truth
dawn/dusk times for a large number of cameras to assess the accuracy
in our estimation of the day/night signal from the image sequences.
This knowledge will help in two ways: it will enable us to quantify
the accuracy of our day/night signal estimate and possibly provide
insight into how to improve the mapping.

\subsubsection{Open Issues}

\begin{itemize}

  \item How much will it help to learn to predict dawn/dusk time for a
    given scene/camera type?

  \item Is conditional correlation the correct similarity metric?  How
    about mutual information?

  \item What distribution is suitable for representing the error in
    localization?

  \item (an interesting question) How unique is the time-series of
    weather conditions for a geographic location?

\end{itemize}


\subsection{Extracting Scene Depth}
\label{sec:depth}

Our goal is to create an automatic method for creating a scene depth
map from a sequence of images captured from a static camera.  The
output will consist of a single image in which the pixel values depend
on the distance from the camera.  The primary cue we will investigate
will be based on geographic consistency of transient weather effects.
This cue is unexplored in the literature and has significant overlap
with our similarity-based camera localization method.

Recent work on depth estimation for outdoor scenes include techniques
using optical effects due to particulate matter in the optical
path~\cite{contrastRestoration, novelDepth2009} and strong structural
priors~\cite{closingTheLoop}. Our approach is more similar to work on
depth-estimation using structured light
patterns~\cite{VaqueroCVPR2009,otherProjectedLight}.  These approaches
require a tight coupling between the camera and the structured light
source.  Our approach differs in only requiring a stochastic
description of the aggregate temporal properties of the light pattern.  

\subsubsection{Motivation}

A depth map is useful for important computer vision applications such
as object detection and object tracking.  In the context of using
webcams as general-purpose sensors, a depth map will likely contribute
useful features for our regression models (\eg to reduce the relative
weights of objects in the distance in a regression model).  In
addition to these practical benefits, exploring the relationship
between the temporal variation of latent factors and the scene depth
will provide valuable insight into the image formation process.

\subsubsection{Approach}

\begin{figure}
\centering

\includegraphics{mountains}

\caption{An ideal afternoon for estimating distance.}

\label{fig:distance}

\end{figure}

In estimating depth, \ie the distance from the camera to the imaged
static scene element, we will use signal correlations.  Previously we
used correlations between images to localize cameras; now we will use
correlations between pixel to estimate the depth.  Two issues must be
addressed: factoring out correlations due to other causes and mapping
the pixel-to-pixel correlation estimates to scene-to-camera distance
estimates.

We are interested in correlations that result from illumination
changes caused by passing clouds.  Considering our geo-temporal image
formation model, we will construct image sets for which this cue is
likely to be the dominant cause of correlation between pixels.  We
will eliminate or reduce the impact of image variations due to: sun
direction, seasonal variation, and human motion.  An ideal image set
will be a few hours of temporally low-pass filtered video with many
small fast-moving clouds (see \figref{distance}).  

In the camera localization algorithm we showed that weather effects,
caused by clouds, have spatial correlations suitable for localization.
Given that sunlight rays are approximately parallel, an image of the
clouds is projected to the ground, much like a satellite image.  If we
ignore scene elements that are not directly illuminated by the sun,
this makes the spatial correlation of light reflected off scene
elements similar to that of the clouds themselves.  

Assuming sufficient samples and an appropriate spatial correlation
function for the clouds, correlations of the intensity of sunlit scene
elements is a one-to-one function of distance between signal sources.
Our goal is to learn this function for the set of images and map
pixel-to-pixel correlations to scene-to-scene distances.  We assume
that there is a single unknown fixed spatial correlation function for
the entire scene and solve for this.

\todo{this is new text\ldots more mathy than the above text\ldots
make it fit together}

Given a sequence of images, each with $N$ pixels, we compute a
similarity score, $s_{i,j}$, for each pairs of pixels. We then
estimate a mapping (see open issues below) that converts the
similarity scores into distances in world coordinates, $d_{i,j} =
f(s_{i,j})$.  A \naive approach to estimating world coordinates from
this distance matrix would be to use multi-dimensional scaling (MDS).
However, MDS does not use the additional constraints provided by the
known camera geometry and can therefore create highly improbable depth
maps. The remainder of this section describes our method for including
camera calibration in the process.

We assume the camera calibration is in the form of a matrix, $U \in
\R^{N \times 3}$, where each row represents the unit ray, in world
coordinates, that corresponds to a pixel.  We solve for the depth map,
$x^*\in\R^{N\times1}$, by estimating a world position for each pixel
subject to the constraint that it lies along the rays specified by the
camera calibration.  This optimization can be conveniently specified
using an inner-product representation as follows:
\begin{equation}
  x^* = \argmin_x \| R(x)  - M(x) \|_F.
\end{equation}
This formulation minimizes the difference in two inner product
matrices, $R$ and $M$. The ray-based matrix,
\begin{equation}
  R(x) = \diag{x}U\trans{U}\diag{x},
\end{equation}
is the the gram matrix of the world points that correspond to the
known depth map and the known pixel rays.  The image measurement-based
gram matrix, $M$, is defined element-wise as:
\begin{equation}
  M_{i,j}(x) = \frac{1}{2}\left( x_i^2 + x_j^2 + d_{i,j}^2 \right).
\end{equation}
This is a the standard transformation between a distance matrix and a
gram matrix with the assumption that the camera center is the origin.

We minimize the objective function using a gradient descent method
(\textbf{fminunc} in Matlab).  The formulation has a straightforward
analytical gradient and Hessian computation.  As an initial condition
we use \todo{something big}.

\subsubsection{Open Issues}
\begin{itemize}
    
  \item What is good representation for the spatial correlation
    function?

  \item How do we convert similarities to distances?  \todo{show some
    pictures of this}

  \item How do we include priors on scene structure to enable
    interpolation across shadow regions?

  \item Is there a better way to translate scene-to-scene distances to
    camera-to-scene distances?

  \item How do we scale this to mega-pixel images?
    
\end{itemize}

\subsubsection{Timeline}

We have an example video that gives promising pixel-to-pixel
correlations.  We have also implemented the embedding 
approach on simulated data.  Depth estimation results for a simulated
$10 \times 10$ image are shown in \figref{depthPreliminary}.  The
preliminary results show high accuracy can be achieved. 

\begin{figure*}
  \centering
  \includegraphics[width=.9\textwidth]{depth_estimation}
  \caption{(left) A ground truth depth map. (right) An estimated depth
  map. \todo{more here} (right) \todo{more here}
  }
  \label{fig:depthPreliminary}
\end{figure*}


\subsection{Learning a Seasonally-Invariant Interest Operator}

NOTE: There are probably already interest operator learning algorithms that
we could just apply to the webcam data.

A two-stage algorithm:

\begin{enumerate}

  \item 
    \begin{itemize}
      \item for each patch build a classifier that separates the patch
        from all others
      \item evaluate the error in the prediction on unseen data (new
        images of the same scene)
    \end{itemize}

  \item
    \begin{itemize}
      \item build a regression model that predicts the error from the
        previous stage given local image features
      \item the regression model is the interest operator
    \end{itemize}

\end{enumerate}

Thoughts:
\begin{itemize}

  \item could be useful for alignment of Internet images

\end{itemize}

\subsection{Controlling for Confounding Geo-temporal Signals}
\label{sec:signalsControlling}

Problem: When we use unsupervised techniques to extract signals we
often get signals that are not interesting.  

Current Approach: Just pick a subset of the data that does not vary
w.r.t. the signals we don't care about and hope for the best.

Solution: Control for known things, day cycle, season cycle, known weather
conditions using some better technique.

We will develop methods for extracting meaningful low-dimensional
time-series signals from webcam images.  An important component will
be reducing the impact of variations due to known sources of
variation.  We will investigate different techniques for reducing the
impact of these nuisance factors.

In addition, some types of image variations, such as motion, will require
different techniques.  We will investigate different approaches to
dealing with motion. 

\subsubsection{Approach}

\todo{this needs to be made manifold learning friendly}

We plan to continue using this approach and to investigate
alternatives for cases where it fails.   Empirical evidence shows
that the top-k PCA method will have difficulty when the underlying
signal of interest is not a dominant cause of variation (the signal
will not be in the basis of the top-k) and when
the related image-space variation is not well explained by a linear
variation (\eg motion).

To handle the case where the underlying signal of interest is not the
dominant source of variation we will remove the impact of known
factors.  Our first approach will be to use a linear technique that
assumes a known, but non-interesting, underlying signal.  We will use
linear regression from this known signal to predict the time-series of
the pixel values.  The residuals of this prediction will be the
variations that are unexplained by the known signal.  We will then use
the standard technique on the residuals to look for the unexplained
sources of variation.

In some cases the confounding time-series may be unknown.  In this
case other techniques may work: solving for the confounding
time-series by selecting pixels that are known to not be affected by
the time-series of interest, imposing constraints on the form of the
variations (\eg must be an affine motion), or pruning the dataset to
remove the non-interesting variations.

In the case of motion, the PCA technique does not give a single
coefficient trajectory that is highly correlated with the underlying
signal of interest.  As a concrete example, consider a set of images
of a circle sliding across a fixed background.   The position of the
circle is the only degree of freedom.  The PCA decomposition of this
image set gives image components that are very similar to a discrete
cosine transform.  The top PCA coefficient trajectory does not
correspond to the position circle.

\subsubsection{Timeline}

We are collecting an initial set of example signal-camera pairs we
plan to use for this investigation.  The first video we have found
contains a flag.  We plan to estimate the distribution of wind
direction given the pixel values.

%
%
%
\subsection{Automatic scene alignment}
\label{sec:sceneAlignment}

\begin{quote}
  This is a very practical contribution but we don't have any ideas
  about how to do it right (other than doing it multi-scale).  It is
  not very exciting.
\end{quote}

Alignment of the {\it semi-static} intervals would more than double the
number of images in the AMOS dataset available to algorithms that
require a {\it static} camera.  There is a long history of
work~\cite{szeliski2006alignment} on joint alignment of image sets and
many of these techniques would work well for alignment of a small set
of webcam images. This section describes a simple algorithm for scene
alignment and highlights the benefits of scene alignment.

Our scene alignment algorithm uses the gradient magnitude image to
reduce sensitivity to weather and illumination conditions.  The
algorithm initializes each image transformation to the identity and
iterates over the following steps until convergence:
    \vspace{-.1in}
\begin{enumerate}
  \item compute the average gradient magnitude image of the
    transformed images,
    \vspace{-.1in}
  \item align each gradient magnitude image to the average gradient
    magnitude image using the Lucas-Kanade method with an affine
    motion model. 
\end{enumerate}
    \vspace{-.1in}
This simple scene alignment algorithm could be improved in many ways:
using better image features for alignment, using a more robust image
distance measure (such as mutual information), automatic support for
large camera motions, and coarse-to-fine alignment.  Developing a
scene alignment technique capable of aligning all the images in the
AMOS dataset with reasonable computational requirements is an
important area for future work.  

\begin{figure}
\centering
\begin{minipage}[b]{0.25\columnwidth} % A minipage that covers half a column
\subfloat[]{ \includegraphics[width=.4\linewidth]{figs/cam_33_stripe_xy} }
\end{minipage}
\hspace{0.05\columnwidth}
\begin{minipage}[ht]{0.25\columnwidth} % A minipage that covers half a column
\subfloat[x-t slice (unaligned)]{ \includegraphics[width=.4\linewidth]{figs/cam_33_stripe_old} }
\vspace{0.05in}
\subfloat[x-t slice (aligned)]{ \includegraphics[width=.4\linewidth]{figs/cam_33_stripe_new} }
\end{minipage}

\caption{Results of the alignment procedure.  (a) An image from the
scene with a horizontal line that shows the image location of the x-t
slices shown below.  (b) An x-t slice from the unaligned image
sequence.  (c) An x-t slice from the aligned image sequence.  }

\label{fig:remove_jitter}
\end{figure}

We qualitatively evaluated the algorithm on several {\it semi-static}
intervals from the AMOS dataset.  Figure~\ref{fig:remove_jitter} shows
results of the scene alignment algorithm on one scene.  To better
understand the impact of this subtle realignment we compared the
principal components of the aligned and unaligned sequences.  The
significant differences in both the components and the coefficients
between the aligned and unaligned sequences are shown in
Figure~\ref{fig:old_new_pca}).  For this scene, in the unaligned
sequence the 3rd PCA component codes seems to code exclusively for
camera motion.  Removing this effect directly impacts PCA based
algorithms for geo-location~\cite{jacobs07geolocate}, and may make
algorithms that infer lighting variation by looking at single pixel
locations over time robust enough to not require hand chosen points.

\begin{figure}
\centering
\subfloat[PCA coefficient scatter plots]{
\includegraphics[width=.4\linewidth]{figs/cam_33_scatters}
}
\subfloat[PCA components]{
\includegraphics[width=.15\linewidth]{figs/cam_33_pca_old}
\hspace{.1in}
\includegraphics[width=.15\linewidth]{figs/cam_33_pca_new}
}

\caption{ (a) Scatter plots of the top three PCA coefficients for the
aligned and unaligned sequences.  The top two coefficients are highly
correlated.  The impact of the alignment step is clearly evident in
the third coefficient (right).  (b) The top three principal component
images for the unaligned (left) and aligned (right) sequence.  The
third component of the unaligned sequence (bottom-left) is
significantly different from that of the aligned sequence.}

\label{fig:old_new_pca}

\end{figure}

\subsection{Other Open Questions}
\label{sec:misc}

This section contains a collection of other tasks that may be
completed as a component of this dissertation.  Some of these items
potentially support some of the major objectives while others extend
them in new directions.

\begin{itemize}

  \item What is the space of object detectors in many scenes?

  \item Can we build a catalog of objects in one scene?

  \item How should you represent images differently depending on
    location (geographic dictionary learning)

  \item How do we find and use motion?  Can we find it using PCA?

  \item What is the distribution of common webcam motions?

  \item Is there a fast (approximate) way to compute mutual
    information between all pairs of patches?

  \item How can we integrate flickr images?

  \item where are all the sources of light in the scene?  can we find
    them during the day and the night?  do they co-vary?

\end{itemize}


